\documentclass[conference]{IEEEtran}
\usepackage{mwe}
\usepackage{cite}
\usepackage{hyperref}

\title{Environmental Impact of Modern Technologies: A Survey}

\author{%
  \IEEEauthorblockN{Jane Doe}
  \IEEEauthorblockA{Department of Environmental Science\\
    University of Example\\
    example@university.edu}
  \and
  \IEEEauthorblockN{Bob Example}
  \IEEEauthorblockA{Institute of Technology\\
    Tech University\\
    bob@techuniversity.edu}
}

\begin{document}

\maketitle

\begin{abstract}
This paper surveys the environmental impact of modern technologies,
focusing on sustainable agriculture, climate change monitoring, and
renewable energy systems. We review recent advances and discuss future
research directions.
\end{abstract}

\section{Introduction}

Modern technologies have a profound impact on the environment. This paper
examines three key areas: sustainable agriculture, climate change
monitoring, and renewable energy systems. Understanding these areas is
essential for developing strategies to mitigate negative environmental
impacts.

\blindtext

\section{Related Work}

Sustainable agriculture has been extensively studied in recent years.
Smith and Brown conducted a comprehensive analysis of sustainable
agriculture methods for tropical regions, demonstrating that organic
farming practices can increase yield by up to 30\% while reducing
chemical input by 50\% \cite{Smith2021}.
Their work identified key factors for successful implementation in
developing countries, including soil quality management and water
conservation techniques.

Machine learning has emerged as a powerful tool for environmental
monitoring. Garcia et al.\ explored machine learning applications in
climate change monitoring, presenting a novel deep learning framework
capable of predicting regional temperature anomalies with 95\%
accuracy \cite{Garcia2022}.
The study analyzed satellite imagery data spanning two decades to train
models that can identify deforestation patterns and carbon emission sources.

Renewable energy systems have undergone rapid development in the past
decade. Chen and Muller provided a comprehensive survey of renewable
energy systems, covering solar, wind, and hydroelectric power
generation \cite{Chen2023}.
They analyzed cost trends, efficiency improvements, and integration
challenges for grid-scale deployment, concluding that solar photovoltaic
systems offer the best cost-to-energy ratio for tropical regions.

\section{Methodology}

\blindtext[2]

\section{Results}

\blindtext[2]

\section{Conclusion}

This survey has reviewed the environmental impact of modern technologies
across three key domains. The work of \cite{Smith2021}, \cite{Garcia2022},
and \cite{Chen2023} collectively demonstrates significant progress in
sustainable technology development.

\bibliographystyle{IEEEtran}
\bibliography{references}

\end{document}
